\documentclass{article}
\usepackage{amsmath, amssymb, amsthm}
\usepackage{hyperref}
\usepackage{geometry}
\geometry{margin=1in}

\title{Erd\H{o}s Problem \#81}
\date{Last updated: 2025-08-31}
\author{Source: erdosproblems.com}

\begin{document}
\maketitle

\noindent\textbf{Status:} OPEN \quad \textbf{No prize}

\noindent\textbf{Tags:} graph theory

\medskip

\noindent\textbf{Problem Statement:}

Let $G$ be a chordal graph on $n$ vertices - that is, $G$ has no induced cycles of length greater than $3$. Can the edges of $G$ be partitioned into $n^2/6+O(n)$ many cliques?




Asked by Erd\H{o}s, Ordman, and Zalcstein \cite{EOZ93}, who proved an upper bound of $(1/4-\epsilon)n^2$ many cliques (for some very small $\epsilon&gt;0$). The example of all edges between a complete graph on $n/3$ vertices and an empty graph on $2n/3$ vertices show that $n^2/6+O(n)$ is sometimes necessary. A split graph is one where the vertices can be split into a clique and an independent set. Every split graph is chordal. Chen, Erd\H{o}s, and Ordman \cite{CEO94} have shown that any split graph can be partitioned into $\frac{3}{16}n^2+O(n)$ many cliques. See also [1017] .




References

[CEO94] Chen, Guan-Tao and Erd\H{o}s, Paul and Ordman, Edward T., Clique partitions of split graphs . Combinatorics, graph theory, algorithms and applications (Beijing, 1993) (1994), 21-30.

[EOZ93] Erd\H{o}s, Paul and Ordman, Edward T. and Zalcstein, Yechezkel, Clique partitions of chordal graphs . Combin. Probab. Comput. (1993), 409-415.

\noindent\textbf{Related OEIS sequences:} possible


\bigskip
\noindent\small{Source: \url{https://www.erdosproblems.com/81}}
\end{document}
